\chapter{HRBR motivation \& design space}
\label{ch:hrbr-motivation}

HRBR (\textbf{H}ybrid \textbf{R}untime with \textbf{B}locks and \textbf{R}eactivity) is the unique differentiator in this codebase.
This chapter exists to justify it as an architecture choice, not as a curiosity.

\section{The problem HRBR is trying to solve}

In the general case, UI rendering has two competing needs:

\begin{itemize}
  \item \textbf{Flexibility:} allow dynamic structure---lists, conditionals, composition, and arbitrary components.
  \item \textbf{Efficiency:} keep updates close to the DOM writes that actually change.
\end{itemize}

Classic Virtual DOM (VDOM) rendering primarily optimizes for flexibility:
we rerender a tree, diff it, then patch the DOM.
That model is incredibly productive.
But it has a recurring cost: \emph{interpretation work} (walk/diff/discover write locations) even when the subtree structure is stable.

HRBR is motivated by a very specific observation:

\begin{quote}
Many real UIs have large subtrees whose \emph{structure} is stable across updates; only values change.
\end{quote}

When structure is stable, paying discovery cost on every update is overhead.

\section{A crisp goal (what would success look like?)}

For stable structure, we want updates to trend toward:

\begin{itemize}
  \item \textbf{Work proportional to changed values} (text/attr/prop/class/style/event), not proportional to subtree size.
\end{itemize}

This is a design goal, not a promise that every benchmark will be faster.
It tells you what HRBR buys when it applies.

\section{Two execution paths, one semantic result}

Relax.js keeps a VDOM baseline for the general case.
HRBR is an optional fast path when the compiler can prove the DOM shape is stable.

\begin{center}
\begin{tikzpicture}[
  node distance=10mm and 14mm,
  box/.style={draw, rounded corners, align=center, inner sep=6pt, fill=RelaxLight},
  arrow/.style={-Latex, thick}
]
\node[box] (author) {Authoring\\TSX/JSX or \texttt{h()}};
\node[box, below left=of author] (vdom) {VDOM runtime\\diff + patch\\general case};
\node[box, below right=of author] (compile) {Compiler\\(optional)};
\node[box, below=of compile] (hrbr) {HRBR runtime\\template + slots\\fast path};
\node[box, below=of author, yshift=-34mm] (dom) {Real DOM\\(same observable result)};

\draw[arrow] (author) -- node[midway,left]{dynamic structure} (vdom);
\draw[arrow] (author) -- node[midway,right]{stable structure} (compile);
\draw[arrow] (compile) -- (hrbr);
\draw[arrow] (vdom) -- (dom);
\draw[arrow] (hrbr) -- (dom);
\end{tikzpicture}
\end{center}

The key point: this book will treat HRBR as an optimization path with a strict boundary.
If the shape isn't stable, HRBR must not guess; it must route to fallback.

\section{Design space: where HRBR sits relative to React/Vue/Solid}

This book focuses on one codebase, but you need a quick coordinate system.
At a grossly simplified level:

\begin{itemize}
  \item \textbf{React:} VDOM reconciliation; great generality; pays update interpretation cost.
  \item \textbf{Vue:} template compiler; reduces work by compiling templates into efficient update code.
  \item \textbf{Solid:} fine-grained reactivity; avoids rerender traversal by tracking dependencies; still needs concrete DOM write sites.
  \item \textbf{Relax HRBR:} stable templates + explicit slot write locations + deterministic scheduling; routes to fallback when structure is dynamic.
\end{itemize}

If you remember one sentence:

\begin{quote}
HRBR trades compiler complexity for fewer categories of runtime work on stable UI structure.
\end{quote}

\section{What you'll see later in the book}

This chapter is motivation.
The rest of the book makes it concrete:

\begin{itemize}
  \item \textbf{Blocks and slots:} what a compiled block is and how slot writes work (Chapter~\ref{ch:hrbr-blocks-template-slots}).
  \item \textbf{Fallback reconciliation:} what happens when structure is dynamic (Chapter~\ref{ch:hrbr-fallback-reconciler}).
  \item \textbf{Compiler boundary:} how we decide block vs fallback, and what the compiler must output (Chapter~\ref{ch:compiler-architecture}).
  \item \textbf{Benchmarks methodology:} how to measure without fooling yourself (Chapter~\ref{ch:benchmarks-methodology}).
\end{itemize}

By the time you reach the compiler chapters, HRBR should feel inevitable: it is a response to a specific performance-shaped problem, with a safety valve for correctness.
